% Continuation of Corollary 3.4.
\noindent a) For every quasi-regular section $a$ of $\mathfrak g$, there
exists, locally for the \'etale topology, a conjugate of $a$ which is a
section of $\mathfrak h$.

\noindent b) For every Cartan subalgebra $\mathfrak d$ of $\mathfrak g$,
$\mathfrak d$ is, locally for the \'etale topology, conjugate to a
subalgebra of $\mathfrak h$.
\end{corollary}
In particular, when $\mathfrak h$ is itself a Cartan subalgebra of
$\mathfrak g$, one finds:

\begin{corollary}{3.5}\label{XIV.3.5}
Let $G$ be a smooth $S$-group prescheme of finite type, $\mathfrak g$
its Lie algebra, and $\mathfrak d,\mathfrak d'$ two Cartan subalgebras
of $G$. Then $\uTransp_G(\mathfrak d,\mathfrak d')$ is identical to
the strict transporter of $\mathfrak d$ into $\mathfrak d'$, and is
a closed subprescheme of $G$ smooth over $S$, with surjective structural
morphism. Locally for the \'etale topology, $\mathfrak d$ and
$\mathfrak d'$ are conjugate.
% typo?: kept as printed: Cartan subalgebras of G.
\end{corollary}
The fact that the transporter is here identical to the strict
transporter follows trivially from the fact that $\mathfrak d$ and
$\mathfrak d'$ are locally direct summands in $\mathfrak g$ and have
the same rank at each point. Thus 3.5 is a particular case of 3.4.
In particular, if $\mathfrak d=\mathfrak d'$:

\begin{corollary}{3.6}\label{XIV.3.6}
Let $G$ be as in 3.5, and let $\mathfrak d$ be a Cartan subalgebra
of $\mathfrak g$. Then $\Norm_G(\mathfrak d)$ is a closed group
subprescheme of $G$ smooth over $S$, whose Lie algebra is identical
to $\mathfrak d$.
\end{corollary}
Indeed, this last assertion amounts to saying that $\mathfrak d$
is its own normalizer in $\mathfrak d$, which follows at once from
the fact that this is true fiber by fiber.
% typo?: kept as printed: normalizer in d rather than g.

\begin{corollary}{3.7}\label{XIV.3.7}
Let $G,\mathfrak g$ be as in 3.5. Then conditions
$(\mathrm C_2),(\mathrm C'_1),(\mathrm C_1)$ of 2.9 are equivalent;
in other words, if $\mathfrak g$ admits a Cartan subalgebra locally
for the fpqc topology, then every quasi-regular section of
$\mathfrak g$ is regular.
\end{corollary}
Indeed, let $a$ be a quasi-regular section; let us prove that it is
regular. The question being local for the fpqc topology, one may
suppose that $\mathfrak g$ admits a Cartan subalgebra $\mathfrak d$.
By 3.4 a), $a$ is then, locally for the \'etale topology, conjugate
to a section of $\mathfrak d$, which reduces one to the case where
$a$ is a section of $\mathfrak d$, where the conclusion is trivial
from the definition.

\begin{definition}{3.8}\label{XIV.3.8}
Let $G$ be a smooth group prescheme over a prescheme $S$. One calls
a subgroup of type (C) of $G$ a group subprescheme $D$ of $G$, smooth
over $S$, with connected fibers, such that $\mathfrak d=\Lie(D)$ is
a Cartan subalgebra (2.4) of $\mathfrak g=\Lie(G)$, i.e. such that
for every $s\in S$, $D_s$ is a subgroup of type (C) of the algebraic
group $G_s$ (XIII 6.2).
\end{definition}

\begin{theorem}{3.9}\label{XIV.3.9}
Let $G$ be a smooth $S$-group prescheme of finite presentation,
and $\mathfrak g$ its Lie algebra. Then:

\noindent a) The map
\[
 D\mto\mathfrak d=\Lie(D)
\]
establishes a one-to-one correspondence between subgroups of type
(C) of $G$ and Cartan subalgebras of $\mathfrak g$.

\noindent b) If $D$ and $\mathfrak d$ correspond, one has
\[
 \Norm_G(D)=\Norm_G(\mathfrak d),
\]
which is a closed subprescheme of $G$ smooth over $S$, and one has
\[
 D=\Norm_G(D)^0=\Norm_G(\mathfrak d)^0.
\]
\noindent c) Two subgroups of type (C), $D$ and $D'$, of $G$ are
conjugate locally for the \'etale topology.
\end{theorem}
\begin{proof}
Let $D$ be a subgroup of type (C) of $G$, and $\mathfrak d$ its
Lie algebra; then $D\subset\Norm_G(\mathfrak d)$, and by definition
3.8 and 3.6 this is an inclusion of group preschemes smooth over
$S$, inducing an isomorphism on the Lie algebras. Since $D$ has
connected fibers, one therefore has $D=\Norm_G(\mathfrak d)^0$.
Thus the map considered in a) is injective; let us prove it is
surjective. Let then $\mathfrak d$ be a Cartan subalgebra of
$\mathfrak g$; then by 3.6, $\Norm_G(\mathfrak d)=N$ is a closed
group subprescheme of $G$ smooth over $S$, admitting $\mathfrak d$
as Lie algebra. Since $G$ is of finite presentation over $S$,
the same holds for $N$, so (as was noted in XII after 7.3) the
union of the connected components of the fibers of $N$ is the
underlying set of an open group subprescheme $N^\circ$ of $N$,
which is clearly a subgroup of type (C) of $G$ with Lie algebra
$\mathfrak d$. This proves a); the first assertion b) follows
at once, and the formula $D=\Norm_G(\mathfrak d)^0$ has already
been proved. Finally, c) follows from a) and 3.5.
% typo?: kept as printed: union of connected components, with no identity qualifier.
\end{proof}

\begin{corollary}{3.10}\label{XIV.3.10}
Suppose that $\mathfrak g$ admits a Cartan subalgebra locally
for the fpqc topology (or, what amounts to the same thing by
3.9 a), that $G$ admits a subgroup of type (C) locally for the
fpqc topology). Consider the functor
$\mathscr D:(\Sch)^\circ_{/S}\to\Ens$ defined by
$\mathscr D(S')=$ the set of subgroups of type (C) of $G_{S'}$.
Then this functor is representable by a quasi-projective prescheme
smooth over $S$, with connected geometric fibers. When $S$ is
the spectrum of a field $k$, hence $G$ a smooth algebraic group
over $k$, and $\mathfrak g$ admits a regular point (a condition
automatically verified if $k$ is infinite), then $\mathscr D$ is
a rational variety over $k$.
\end{corollary}
Indeed, by 3.9 a), the functor $\mathscr D$ is canonically
isomorphic to the functor considered in 2.16; on the other hand,
by 3.7, condition $(\mathrm C_2)$ holds. Thus 3.10 follows from
2.16 and 2.17.

\begin{corollary}{3.11}\label{XIV.3.11}
Let $D$ be a subgroup of type (C) of $G$, and $N$ its normalizer
in $G$. Then the quotient sheaf $G/N$ is canonically isomorphic
to the functor $\mathscr D$ of 3.10, hence representable by a
quasi-projective prescheme smooth over $S$, with connected
geometric fibers.
\end{corollary}

\begin{proposition}{3.12}\label{XIV.3.12}
Let $G$ be an $S$-group prescheme smooth and of finite
presentation over $S$, $H,K$ two group subpreschemes smooth
and of finite presentation over $S$, $K$ having connected
fibers, and $\mathfrak g,\mathfrak k,\mathfrak h$ the corresponding
Lie algebras. One assumes that one of the following two
conditions holds for the geometric fibers of the latter:

\noindent a) For every $s\in S$, $\mathfrak k_{\bar s}$ contains
a regular element of $\mathfrak g_g$, and $\mathfrak h_{\bar s}$
contains a Cartan subalgebra of $\mathfrak g_{\bar s}$.
% typo?: kept as printed: g_g in a).

\noindent b) For every $s\in S$, $\mathfrak k_{\bar s}$ contains
a Cartan subalgebra of $\mathfrak g_{\bar s}$. Under these
conditions, for $H\supset K$ to hold, it is necessary and
sufficient that one have $\mathfrak h\supset\mathfrak k$.
\end{proposition}
Of course, one need only prove that if
$\mathfrak h\supset\mathfrak k$, then $H\supset K$.
In case b), the inclusion $\mathfrak h\supset\mathfrak k$
shows that one is in fact under the conditions of a),
so it suffices to prove a). Proceeding as in 3.2 a)
by reduction to the case of $S$ artinian local, one
is reduced to the case where there exists a section
$a$ of $\mathfrak k$ which is quasi-regular in
$\mathfrak g$. In this case, proceeding as in XIII
5.5, one is reduced to the following statement:

\begin{corollary}{3.13}\label{XIV.3.13}
Let $G$ be an $S$-group prescheme smooth and of finite
presentation over $S$, $H$ a group subprescheme of
$G$ smooth and of finite presentation over $S$,
$\mathfrak g$ and $\mathfrak h$ the Lie algebras,
and $a$ a section of $\mathfrak h$ which is
quasi-regular in $\mathfrak g$. One assumes that
for every $s\in S$, the geometric fiber
$\mathfrak h_{\bar s}$ contains a Cartan subalgebra
of $\mathfrak g_{\bar s}$. Let
\[
 M_a=\Transp_G(a,\mathfrak h),
\]
which is a closed subprescheme of $G$ smooth over
$S$ (cf. 3.1), so that $M_a^0$ (the union of the
connected components of the identity element in
the fibers of $M_a$) is an open part of $M_a$,
which we shall equip with the structure induced
by $M_a$. One then has $H^0=M_a^0$.
\end{corollary}
Clearly, one has $H\subset M_a$, whence
$H^0\subset M_a^0$. Since this is an inclusion
of preschemes smooth over $S$, to prove that
it is an equality, one is reduced to verifying
it on the fibers, which reduces one to the
case where $S$ is the spectrum of an
algebraically closed field, a case seen in
XIII 5.4.

\begin{corollary}{3.14}\label{XIV.3.14}
Let $G$ be a smooth $S$-group prescheme of
finite presentation, $D$ a subgroup of type
(C) of $G$, and $H$ a group subprescheme of
$G$ smooth over $S$ and of finite presentation
over $S$, and $\mathfrak d$ and $\mathfrak h$
their Lie algebras. Then one has
\[
 \uTransp_G(D,H)=\uTransp_G(\mathfrak d,\mathfrak h),
\]
and this functor is representable by a
closed subprescheme of $G$ smooth over $S$.
\end{corollary}
Indeed, the identity between the two
transporters follows from 3.12 b), which
allows one to apply 3.2.

\begin{corollary}{3.15}\label{XIV.3.15}
Let $G,H$ be as in 3.14, and suppose that
for every $s\in S$, the geometric fiber
$\mathfrak h_{\bar s}$ contains a Cartan
subalgebra of $\mathfrak g_{\bar s}$. Suppose
moreover that $\mathfrak g$ admits a Cartan
subalgebra locally for the fpqc topology.
Then locally for the \'etale topology,
$H$ contains a subgroup of type (C) of $G$.
\end{corollary}
By 3.7 and 3.9 a), $G$ admits a subgroup
of type (C) locally for the \'etale topology,
so one may suppose that $G$ admits such
a subgroup, say $D$. Then the hypothesis
on $\mathfrak h$ also means that the
structural morphism of the transporter
considered in 3.14 is surjective (taking
into account the conjugacy theorem XIII
6.1 a)). One then concludes by Hensel's
lemma XI 1.10.

\begin{corollary}{3.16}\label{XIV.3.16}
Let $G,H,K$ be as in 3.12 a); suppose
moreover that for every $s\in S$, the
geometric fiber $\mathfrak k_{\bar s}$
is nilpotent (i.e. $\mathfrak k$ is
locally nilpotent). Then one has
\[
 \uTransp_G(K,H)=\uTransp_G(\mathfrak k,\mathfrak h),
\]
and this functor is representable by
a closed subprescheme of $G$ smooth
over $S$, with surjective structural
morphism to $S$. Locally for the
\'etale topology, $H$ contains a
subgroup conjugate to $K$.
\end{corollary}
The identity of the two transporters
is again contained in 3.12 a); the
assertion about its structure is
then precisely 3.2 a), and the last
assertion of 3.16 is then a consequence
of Hensel's lemma.

\begin{remarks}{3.17}\label{XIV.3.17}
\noindent a) In 3.12 and 3.16, one
may replace the hypothesis that
$K$ is smooth over $S$ by the
following weaker hypothesis:
the sheaf $e_K^*(\Omega^1_{K/S})=\omega_K^1$
of relative $1$-differentials along
the unit section is locally free.
In this way, 3.12 contains XIII 5.5.

\noindent b) Let $G,\mathfrak g,\mathfrak h$
be as in 3.1, with $G$ of finite
presentation over $S$. Then
$N=\Norm_G(\mathfrak h)$ is not
necessarily smooth over $S$ along
the unit section, or, what amounts
to the same thing, there does not
necessarily exist a group
subprescheme $H$ of $S$ smooth
over $S$ whose Lie algebra is
$\mathfrak h$, even if $S$ is
the spectrum of a field. When
such an $H$ exists, so that one
then has (taking $H$ with
connected fibers) $N=\Norm_G(H)$,
I do not know whether $N$ is
smooth over $S$. In this question,
one may clearly reduce to the
case where $S$ is artinian local.
% typo: il n'existe par -> il n'existe pas.
% typo?: kept as printed: a group subprescheme H of S.
\end{remarks}
Finally, to end this section,
let us examine the case where
$G$ is ``semisimple'' over $S$:

\begin{theorem}{3.18}\label{XIV.3.18}
Let $S$ be a prescheme, $G$ an
$S$-group prescheme smooth over
$S$, whose geometric fibers
are ``adjoint'' semisimple
groups, i.e. semisimple with
reductive center (XII 4.1 and
4.4) reduced to the unit group.
Then the subgroups of type (C)
of $G$ are identical to its
maximal tori, hence also to
its Cartan subgroups (XII 3.1).
\end{theorem}
Taking the definitions into
account, one is reduced to the
case where $S$ is the spectrum
of an algebraically closed
field, and to proving then
that for a maximal torus $T$
of $G$, the Lie algebra
$\mathfrak t$ of $T$ is a
Cartan subalgebra of
$\mathfrak g$, that is
(taking into account the
inequality: nilpotent rank
of $\mathfrak g\geq$ nilpotent
rank of $G=\dim T=
\operatorname{rank}\mathfrak t=r$)
that there exists
$x\in\mathfrak t$ with
$\dim\operatorname{Nil}(x,\mathfrak g)=r$.
Since $\mathfrak t$ is
abelian and a fortiori
nilpotent, it amounts to
the same thing to say that
there exists an
$a\in\mathfrak t$ such
that $\ad(a)_{\mathfrak g/\mathfrak t}$
is injective (XIII 5.7 a)).
Now consider the characters
$\alpha$ of $T$ which occur
in the representation induced
on $T$ by the adjoint
representation of $G$.
The structure theory of the
semisimple group $G$ (BIBLE,
13 th. 1 a) and th. 3,
cor. 2), more precisely of
the ``big cell'' of $G$,
a semidirect product of
$T$ and subgroups $P_\alpha$
isomorphic to the additive
group $\Ga$, invariant
under $T$ and corresponding
to the ``root'' characters
of $G$ for the torus $T$,
shows that the eigenspace
of $\mathfrak g$ for the
unit character is precisely
$\mathfrak t$, and the
other eigenspaces have
dimension $1$, the associated
characters $\alpha$ being
precisely the roots of $G$
for $T$. By the calculation
of the reductive center
of $G$ as the intersection
of the kernels of the
characters of $T$ which
occur in the adjoint
representation of $G$
(XII 4.8), one sees that
$G$ being adjoint is
interpreted by the fact
that the roots generate
the lattice
$M=\Hom(T,\mathbf G)$.
Now a well-known lemma
of root theory tells
us that every root
belongs to a system
of simple roots, hence
to a basis of the
group generated by the
roots, and consequently
to a basis of the
dual $M$ of $T$\nde{$\widehat T$
has been replaced by $T$.}.
One concludes:
% typo?: kept as printed: Hom(T,G) with no m subscript.

\begin{corollary}{3.19}\label{XIV.3.19}
If $G$ is an adjoint
semisimple algebraic group
over an algebraically
closed field $k$, and
$T$ a maximal torus
of $G$, then for every
root $\alpha$ of $G$
with respect to $T$,
$\alpha:T\to\Gm$, the
corresponding homomorphism
$\alpha':\mathfrak t\to k$
is nonzero.
\end{corollary}
This result is essentially
equivalent to theorem
3.18, for, for a
$t\in\mathfrak t$,
$\ad(t)$ is semisimple,
and its eigenvalues
on $\mathfrak g/\mathfrak t$
are precisely the
$\alpha'(t)$; hence
$\ad(t)_{\mathfrak g/\mathfrak t}$
is injective if and
only if the $\alpha'(t)$
are $\ne0$, and there
exists a $t\in\mathfrak t$
with this property if
and only if all the
$\alpha'$ are $\ne0$.

\begin{corollary}{3.20}\label{XIV.3.20}
Let $S$ be a prescheme,
$G$ an $S$-group prescheme,
smooth and of finite
presentation over $S$,
with geometric fibers
connected reductive
algebraic groups (i.e.
extensions of a
semisimple group by
a torus). Then for
every $s\in S$ there
exists an open
neighborhood $U$ of
$s$ such that $G|_U$
admits a maximal torus.%
{\renewcommand{\thefootnote}{*}\footnote{This
result, as well as
2.11 on which it
rests, generalizes
immediately to the
case where $s$ is
replaced by a finite
part of $S$ contained
in an affine open
subset.}\addtocounter{footnote}{-1}}
\end{corollary}
Indeed, we shall see
in XVI that the
hypothesis just made
on $G$ implies that
$G$ is affine over
$S$ and has locally
constant reductive
rank; hence (XII
4.7 c)) $G$ admits
a reductive center
$Z$, $G'=G/Z$ is
a smooth affine
group over $S$ whose
reductive center is
reduced to the unit
group, and finally
the maximal tori
of $G$ and of $G'$
are in one-to-one
correspondence. Moreover,
one sees at once
that the geometric
fibers of $G'$ are
connected semisimple
groups, and moreover
adjoint by definition
(their reductive center
being trivial). Thus
it suffices to confine
oneself to the case
where $G$ is semisimple
and adjoint, and by
3.18 one is reduced
to finding an open
neighborhood $U$ of
$s$ and a subgroup
of type (C) of $G|_U$,
or, what amounts to
the same thing
(3.9 a)), a Cartan
subalgebra of
$\mathfrak g|_U$.
Now this is possible
if $k(s)$ is infinite,
for by 3.7,
$\mathfrak g$ satisfies
condition $(\mathrm C_2)$
of 2.9, hence one
may apply 2.11 b).
In fact, statement
3.20 remains valid
without supposing
$k(s)$ infinite.
Indeed, by the
preceding proof,
it suffices to know
that for every
adjoint semisimple
group $G$ over a
finite field $k$,
the Lie algebra
$\mathfrak g$ of
$G$ contains a
regular element.
Now this statement
has been proved by
Chevalley (using the
properties of the
Coxeter element
of the Weyl group\ldots);
cf. the Appendix
below by J.-P. Serre.

\begin{remarks}{3.21}\label{XIV.3.21}
\noindent a) Statement
3.20 remains valid,
with essentially the
same proof, replacing
$G$ by a closed group
subprescheme $H$,
smooth over $S$,
having everywhere
the same rank as $G$
(for example, a
``parabolic subgroup''
of $G$), provided
$k(s)$ is infinite.
I do not know whether
here again the
hypothesis that
$k(s)$ be infinite
is superfluous.
One can show that
it is so at least
if $H$ is parabolic,
thanks to the
construction of the
radical $U$ of $H$
and of the quotient
$H/U$, which reduces
one to the semisimple
case. Unfortunately,
the method of regular
elements seems powerless
here, for one easily
constructs examples
(for example, with
the projective group
and its ``standard''
Borel subgroup) where
the Lie algebra
$\mathfrak h$ of $H$
contains no regular
element.

\noindent b) The
proof of 3.20 in
fact shows a more
precise result
(invoking 3.9 b))
in the case where
$k(s)$ is infinite,
namely that every
maximal torus $T_0$
of $G_s$ comes from
a maximal torus $T$
over an open
neighborhood of $s$.
I do not know whether
this statement remains
valid when $k(s)$
is no longer supposed
infinite, the difficulty
clearly coming from
the fact that the
Lie algebra
$\mathfrak t_0$ of
$T_0$ generally
contains no regular
element of the
Lie algebra
$\mathfrak g_0$ of
the fiber $G_s$.
An affirmative answer
to this problem would
imply the following
statement (which is
proved only in the
case of an infinite
residue field or when
$A$ is separated and
complete): Let $A$
be a local ring
with residue field
$k$, $M$ an ``Azumaya
algebra'' over $A$,
i.e. an algebra
such that $M$ is
a free module of
finite type over $A$
and $M_0=M\otimes_A k$
a central simple
algebra over $k$,
and $D_0$ a
commutative subalgebra
of $M_0$ separable
over $k$, such that
$[M_0:k]=([D_0:k])^2$;
then there exists
a commutative
subalgebra $D$ of
$M$ which is a
module direct summand
in $M$ and such
that $D\otimes_A k=D_0$
(?). (Note that
giving $M$ is
equivalent to giving
a principal homogeneous
bundle under the
projective group
$\PGL(n)_A$, whence
an ``inner'' twisted
form $G$ of $\PGL(n)$,
and the maximal tori
of $G$ correspond
one-to-one to the
commutative subalgebras
$D$ of $M$, \'etale
over $A$, of rank
$n$ over $A$.)

\noindent c) Applying
3.20 to the
centralizer of a
subtorus $Q$ of $G$
($G$ a reductive
$S$-group prescheme),
one deduces that
every such $Q$ is
contained, locally
for the Zariski
topology, in a
maximal torus of $G$.
\end{remarks}

\sgasection{4}{A digression on Borel subgroups}
\begin{definition}{4.1}\label{XIV.4.1}
Let $G$ be a smooth
algebraic group over
an algebraically
closed field. One
calls a Borel subgroup
of $G$ a smooth
solvable connected
algebraic subgroup
which is maximal
for these properties.
\end{definition}
When $G$ is affine,
one thus recovers
the terminology of
BIBLE 6 def. 1.
Note immediately
that if $Z$ is a
smooth connected
subgroup of $G$
contained in the
center (or, more
generally, solvable
and invariant), then
for every Borel
subgroup $B$ of $G$,
the image $BZ$ of
$B\times Z$ under
the morphism
$(b,z)\mto bz$
from $B\times Z$
to $G$ is a
smooth solvable
connected subgroup
of $G$ containing
$B$, hence identical
to $B$, so $B$
contains $Z$, hence
$B$ is the inverse
image of an
algebraic subgroup
$B'$ of $G'=G/Z$,
and it is immediate
that $B'$ is a
Borel subgroup of
$G'$. Taking
$Z=\Centr(G^0)^0_{\mathrm{red}}$,
$G'=G/Z$ is affine
(XII 6.1), so the
Borel subgroups of
$G'$ are conjugate,
and for such a
$B'$, $G'/B'$
is a projective
variety (BIBLE 6
th. 4). Consequently:

\begin{proposition}{4.2}\label{XIV.4.2}
Let $G$ be as in
4.1. Then the Borel
subgroups of $G$
are conjugate.
If $B$ is a Borel
subgroup, then
$G/B$ is a
projective variety.
The maximal tori
of $B$ (resp. the
Cartan subgroups of
$B$, $G$ being
connected) are
maximal tori of
$G$ (resp. Cartan
subgroups of $G$).
\end{proposition}
